\chapter{Алгоритмическая оптимизация}
\label{ch:algorithmic}

\section{Общий буфер окна и упреждающего буфера}

\textbf{Проблема.} Исходная программа хранит окно (4~КиБ уже сжатых данных)
в кольцевом буфере, а упреждающий буфер "--- в отдельном массиве того же
размера. После каждого токена весь остаток упреждающего буфера
сдвигается к его началу (листинг~\ref{lst:memmove-before}). Как показано в
разделе~\ref{sec:base-complexity}, это член $O(T \cdot L)$ в оценке
сложности: на несжимаемых данных копируется около 2~КБ на каждый входной
байт. Кроме того, каждый байт кандидата в цикле сравнения читается
функцией \texttt{byte\_at}, которая с ветвлением выбирает между двумя
массивами и вычисляет остаток от деления.

\begin{lstlisting}[caption={Сдвиг упреждающего буфера после каждого токена (исходная программа)},label={lst:memmove-before}]
while (len < MAX_MATCH && len < buf_len
       && byte_at (window, buf, win_pos, pos + len) == buf[len])
  len++;
/* ... emit the token ... */
memmove (buf, buf + advance, buf_len - advance);
buf_len -= advance;
if (buf_len < MAX_MATCH)
  buf_len += fread (buf + buf_len, 1, sizeof (buf) - buf_len, fin);
\end{lstlisting}

\textbf{Решение.} Окно и упреждающий буфер объединены в один массив
\texttt{data} размером 64~КиБ (листинг~\ref{lst:memmove-after}).
Переменная \texttt{data\_base} хранит абсолютную позицию первого байта
массива в файле, а текущая позиция "--- это просто указатель
\texttt{cursor} внутри массива. Переход к следующему токену сводится к
увеличению счётчика позиции. Когда в массиве остаётся меньше 4~КиБ
непрочитанных данных, он дополняется одним вызовом \texttt{fread}; если места
не хватает, последние 4~КиБ истории и остаток впереди (не более 8~КиБ)
переносятся в начало массива. Такой перенос происходит один раз примерно
на 56~КиБ входных данных вместо одного сдвига на каждый токен.

Кандидат совпадения и упреждающий буфер теперь "--- два указателя в один
массив, поэтому функция \texttt{byte\_at} больше не нужна, а цикл
сравнения стал простым сравнением двух последовательностей байт. Длина
упреждающего буфера \texttt{look\_len} вычисляется так же, как в исходной
программе, потому что от неё зависит, какие позиции в конце блока попадают
в хеш-таблицу: это сохраняет архивы неизменными.

\begin{lstlisting}[caption={Дополнение общего буфера (после оптимизации)},label={lst:memmove-after}]
look_len -= advance;
if (look_len < MAX_MATCH)
  {
    size_t start = (size_t) (win_pos - data_base);
    size_t ahead = data_len - start;
    if (ahead < WINDOW_SIZE)
      {
        if (data_len + (WINDOW_SIZE - ahead) > sizeof (data))
          {
            size_t keep_from = start - WINDOW_SIZE;
            memmove (data, data + keep_from, data_len - keep_from);
            data_len -= keep_from;
            data_base += keep_from;
            start -= keep_from;
          }
        data_len += fread (data + data_len, 1,
                           sizeof (data) - data_len, fin);
        ahead = data_len - start;
      }
    look_len = ahead < WINDOW_SIZE ? ahead : WINDOW_SIZE;
  }
\end{lstlisting}

\textbf{Изменение сложности.} Объём копирования уменьшается с
$O(T \cdot L)$ до $O(n)$: каждый байт переносится не чаще одного раза на
блок, а доля перенесённых байт не превышает $8/56 \approx 0{,}15$ от объёма
входа. Оценка~\eqref{eq:base-compress} принимает вид
\begin{equation}
\label{eq:after-14}
T_{\text{сж}}(n) = O\bigl(n + T \cdot C \cdot M\bigr),
\end{equation}
то есть из стоимости каждого токена исключено слагаемое $L = 4096$.
Поиск совпадений при этом не меняется.

\textbf{Корректность.} Проверены \texttt{make check} и совпадение архивов
с архивами исходной программы для всех файлов корпуса, для граничных
случаев (пустой файл, файлы длиной 1--9, 255, 256, 4095--4097,
65\,535--65\,537 байт) и для фрагментов по 1,5~МБ, пересекающих многие
переносы буфера.

\section{Результаты алгоритмической оптимизации}

Результаты приведены в таблице~\ref{tab:res-14}. Сжатие всего корпуса стало
в 1,52 раза быстрее на debian и в 1,29 раза на hppavillion; на файлах
корпуса Silesia "--- в 1,58 и 1,51 раза. Распаковку эта оптимизация не
затрагивает.

\begin{table}[H]
\caption{Результаты алгоритмической оптимизации (шаг 03 против шага 02)}
\label{tab:res-14}
\begin{tabletext}
\centering
\begin{tabular}{|l|r|r|r|r|r|r|}
\hline
\multirow{2}{*}{\textbf{Метрика}} & \multicolumn{3}{c|}{\textbf{debian}} & \multicolumn{3}{c|}{\textbf{hppavillion}} \\
\cline{2-7}
 & \textbf{до} & \textbf{после} & \textbf{ускор.} & \textbf{до} & \textbf{после} & \textbf{ускор.} \\
\hline
Сжатие корпуса, млрд тактов & 41,94 & 27,62 & \textbf{1,52} & 40,86 & 31,65 & \textbf{1,29} \\
\quad только Silesia & 22,72 & 14,35 & 1,58 & 20,99 & 13,90 & 1,51 \\
Сжатие корпуса, с & 9,89 & 7,04 & 1,41 & "--- & "--- & "--- \\
dickens: такты, млн & 1071 & 785 & 1,36 & 1135 & 784 & 1,45 \\
dickens: команды, млн & 2428 & 1366 & 1,78 & 2406 & 1322 & 1,82 \\
dickens: IPC & 2,27 & 1,74 & & 2,12 & 1,68 & \\
\hline
\end{tabular}
\end{tabletext}
\end{table}

Число команд уменьшилось почти вдвое, а число тактов "--- меньше: удалённое
копирование \texttt{memmove} выполнялось с высоким числом команд за такт
(широкие векторные пересылки), поэтому средний IPC программы упал с 2,27 до
1,74. Теперь время определяется поиском совпадений с его зависимыми
переходами по цепочкам.

Выигрыш на отдельных файлах (рисунок~\ref{fig:per-file-14}) тем больше, чем
больше токенов приходится на байт входа, ведь сдвиг выполнялся один раз на
токен. Файл random.dat, состоящий почти из одних литералов, сжимается в
3,8~раза быстрее; двоичные файлы с короткими совпадениями (osdb, sao,
x-ray, ooffice) "--- в 1,5--2,2~раза; тексты "--- примерно в 1,3~раза;
zeros.dat, где каждый токен "--- совпадение максимальной длины, "---
меньше всего.

\begin{figure}[H]
\centering
\begin{tikzpicture}
\begin{axis}[
  ybar, width=\textwidth, height=7cm, bar width=8pt,
  ylabel={Ускорение сжатия, раз}, ymin=0,
  xtick=data, xticklabels from table={data/perfile-debian.dat}{file},
  xticklabel style={rotate=60, anchor=east, font=\small\ttfamily},
  nodes near coords, nodes near coords style={font=\tiny, rotate=90, anchor=west},
  enlarge x limits=0.04, grid=major, grid style={gray!25}, ymax=4.6,
]
\addplot[fill=gray!60] table[x=index, y=s14] {data/perfile-debian.dat};
\draw[dashed] (axis cs:-1,1) -- (axis cs:16,1);
\end{axis}
\end{tikzpicture}
\caption{Ускорение сжатия по файлам после алгоритмической оптимизации (debian, время)}
\label{fig:per-file-14}
\end{figure}

\screenshot{perf-14}{Метрики perf stat на файле dickens после алгоритмической оптимизации}{fig:perf-14}{%
\texttt{sed -n '/perf stat, compress/,/elapsed/p'
bench/log/debian/03-no-memmove/raw/profile/summary.txt}; рядом то же для
\texttt{02-inline-bitio}, чтобы видеть <<до>> и <<после>>.}
